% MSc BIS Thesis 
% LaTeX Template
% Version 1.0 (29/08/2014)
%
% Original authors:
% Steven Gunn 
% Sunil Patel
%
% Further authors:
% Michael Stauffer
% 
% This template is based on
% http://www.latextemplates.com/template/masters-doctoral-thesis
%
% License:
% CC BY-NC-SA 3.0 (http://creativecommons.org/licenses/by-nc-sa/3.0/)
%
%%%%%%%%%%%%%%%%%%%%%%%%%%%%%%%%%%%%%%%%%

%----------------------------------------------------------------------------------------
%	PACKAGES AND OTHER DOCUMENT CONFIGURATIONS
%----------------------------------------------------------------------------------------

% For digital publication it is recommended to use a one-sided layout
\documentclass[11pt, oneside]{Thesis}

\pgfplotsset{compat=1.18}

\addbibresource{Bibliography.bib}

% For printing it is recommended to use a two-sided layout
%\documentclass[11pt, twoside]{Thesis}

\setlength{\epigraphwidth}{.7\textwidth}
\renewcommand{\epigraphflush}{flushright}
\makeatletter
% Taken and updated from http://mirrors.ctan.org/macros/latex/contrib/epigraph/epigraph.dtx
\renewcommand{\@epitext}[1]{%
  \begin{minipage}{\epigraphwidth}\begin{\textflush} \hspace*{0pt}#1\\
    \ifdim\epigraphrule>\z@ \@epirule \else \vspace*{-.5\baselineskip} \fi
  \end{\textflush}\end{minipage}}
\makeatother


% SPARQL lst sintax
\lstset{breaklines=true}
\definecolor{LightGray}{rgb}{0.97,0.97,0.97}
\lstdefinelanguage{SPARQL}{
  basicstyle=\footnotesize\ttfamily,
  backgroundcolor=\color{LightGray},
  columns=fullflexible,
  breaklines=true,
  sensitive=true,
  % --------------------------
  frame=bt,
  aboveskip=1em,
  belowskip=1em,
  xleftmargin=.5em,
  xrightmargin=.5em,
  framexleftmargin=.5em,
  framextopmargin=.5em,
  framexbottommargin=.5em,
  framexrightmargin=.5em,
  % --------------------------
  tabsize = 2,
  showstringspaces=false,
  morecomment=[l][\color{gray}]{\#},       % comments
  morecomment=[n][\color{blue}]{<http}{>}, % uris
  morestring=[b][\color{OliveGreen}]{\"},  % strings
  % -------------------------- variables
  keywordsprefix=?,
  classoffset=0,
  keywordstyle=\color{Sepia},
  morekeywords={},
  % -------------------------- prefixes
  classoffset=1,
  keywordstyle=\color{Purple},
  morekeywords={rdf,rdfs,owl,xsd,purl, eurio},
  % -------------------------- keywords
  classoffset=2,
  keywordstyle=\color{MidnightBlue},
  morekeywords={
    SELECT,CONSTRUCT,DESCRIBE,ASK,WHERE,FROM,NAMED,PREFIX,BASE,OPTIONAL,
    FILTER,GRAPH,LIMIT,OFFSET,SERVICE,UNION,EXISTS,NOT,BINDINGS,MINUS,a
  }
}

\begin{document}

%----------------------------------------------------------------------------------------
%	DOCUMENT VARIABLES
%	Fill in the lines below to update the thesis template
%	
%   If you want to cite each of the variables defined below, look at the
%	section above for the citation command.
%----------------------------------------------------------------------------------------

% You thesis title
% Citation command \thesisTitle
\thesisTitle{Title to be defined}
%-------------------------------------------------  

% You supervisor's name Unicam or FHNW
% Citation command \supName 
\supervisor{Emanuele Laurenzi}
%-------------------------------------------------  

% You supervisor's name Unicam or FHNW
% Citation command \supName 
\supervisorTwo{Massimo Callisto De Donato}
%-------------------------------------------------  

% If you don't have a co-supervisor, exchange with
%\cosupervisorfalse
\cosupervisortrue
%-------------------------------------------------

% You co-supervisor's name 
% Citation command \supcoName 
\cosupervisor{Piermichele Rosati}
%-------------------------------------------------  

% Your degree name
% Citation command \degreeName
\degree{Business Information Systems}
\degreeTwo{Computer Science}
%-------------------------------------------------   

% Your name
% Citation command \authorName
\authors{Eduardo Tiano}
\matricola{129508 / 25-657-446}%
\academicyear{2025/2026}%
%-------------------------------------------------

%----------------------------------------------------------------------------------------
%	GLOSSARY
%	Add your glossary terms below
%	
%   If you want to cite a glossary term simply use \gls{term) for the singular form 
%   of the term or \glspl{term} for the plural form of the term.
%----------------------------------------------------------------------------------------

\newglossaryentry{Java}{
    name={Java},
    description={\blindtext}
}

\newglossaryentry{Oracle}{
    name={Oracle},
    description={\blindtext}
}

\newglossaryentry{Microsoft}{
    name={Microsoft},
    description={\blindtext}
}

%----------------------------------------------------------------------------------------
%	ABBREVIATIONS
%	Add your glossary terms below
%	
%   To generate the list of abbreviations run the following command in the terminal:
%   makeglossaries Rosati-2024-Thesis 
%
%   If you want to cite an abbreviation term simply use \gls{term) for the singular form  
%   of the term or \glspl{term} for the plural form of the term.
%----------------------------------------------------------------------------------------

\newacronym{rdf}{RDF}{Resource Description Framework}
\newacronym{rdfs}{RDFS}{RDF Schema}
\newacronym{uri}{URI}{Uniform Resource Identifier}
\newacronym{owl}{OWL}{Web Ontology Language}
\newacronym{sparql}{SPARQL}{SPARQL Protocol and RDF Query Language}
\newacronym{nt}{NT}{N-Triples}
\newacronym{ttl}{TTL}{Turtle}
\newacronym{jsonld}{JSON-LD}{JavaScript Object Notation for Linked Data}
\newacronym{csv}{CSV}{Comma-Separated Values}
\newacronym{bim}{BIM}{Building Information Modeling}
\newacronym{ss}{SS}{Answer Semantic Similarity}
\newacronym{nan}{NaN}{Not a Number}

% PDF meta-data
\hypersetup{pdftitle={\thesisTitle}}
\hypersetup{pdfauthor=\authorName}

\frontmatter % Use roman page numbering style (i, ii, iii, iv...) for the pre-content pages

%----------------------------------------------------------------------------------------
%	TITLE PAGE
%----------------------------------------------------------------------------------------

\titlePage
%\clearpage
%\input{chapters/quotation.tex}
%\clearpage

%----------------------------------------------------------------------------------------
%	ABSTRACT PAGE
%----------------------------------------------------------------------------------------
\pagenumbering{roman}
%\input{chapters/Abstract}

\clearpage % Start a new page
%\input{chapters/Abstract-ita}

\clearpage % Start a new page

%----------------------------------------------------------------------------------------
%	STATEMENT OF AUTHENTICITY
%----------------------------------------------------------------------------------------

\fillingPage{} 
%\statementOfAuthenticity

\clearpage

\tableofcontents

\fillingPage{} 
\listOfFigures

\fillingPage{} 
\listOfTables

\mainmatter % Begin numeric (1,2,3...) page numbering

\fillingPage{}
\chapter{Problem Statement}
\label{chap:problemStatement}

Internet of Things (IoT) is a term denoting interconnection of devices connected to a network,
capable of communicating via different means \cite{berte_defining_iot_2018}.  
\fillingPage{}
%\input{chapters/LiteratureReview}

%----------------------------------------------------------------------------------------
%	BIBLIOGRAPHY
%----------------------------------------------------------------------------------------

\fillingPage{}
\nocite{*}
\listOfBibliography

%----------------------------------------------------------------------------------------
%	THESIS CONTENT - APPENDICES (if appropriate)
%----------------------------------------------------------------------------------------

\appendix % Cue to tell LaTeX that the following 'chapters' are Appendices

% Include the appendices of the thesis as separate files from the Appendices folder
% Uncomment the lines as you write the Appendices

\fillingPage{} 
%\input{appendices/ProjectParticipantsAgent}

%----------------------------------------------------------------------------------------
%	FOREWORD / ACKNOWLEDGEMENTS (optional)
%----------------------------------------------------------------------------------------

%\input{chapters/Acknowledgments}

\clearpage % Start a new page

\backmatter

\end{document}  